\normalsize{ \indent
Para el software completo, se proponen los siguientes objetivos
a largo plazo:
}
\begin{itemize}
  \item Poder tener los horarios de los cursos disponibles,
  dentro de la solución final; donde sea posible rotar
  horarios, gestionar licencias, altas y bajas. Es muy
  importante controlar que cada una de estas acciones se
  realice correctamente, siguiendo todas las reglas que tiene
  el colegio respecto a este tema.
  \item Crear una base de datos, donde se tenga los componentes
  necesarios de un legajo docente; con la opción de observar
  éstos de forma ordenada y fácilmente, a la vez tener un acceso
  rápido a la información (sin tener que buscar en una planilla
  de Excel o en carpetas de archivos, ni mucho menos en formularios
  manuscritos).
  \item Tener un apartado dedicado a los documentos históricos que
  existieron antes de la solución propuesta, donde se tendrá un
  ordenamiento eficaz para todo usuario que quiera y tenga permisos
  para acceder a tal información. Inclusive, todos los documentos
  que sean relevados por otros irán a esta sección.
  \item Hallar la forma de digitalizar los documentos posibles;
  así se evita perder tiempo en el paso de datos. Sin embargo,
  cabe destacar que también hay que generar el formulario físico;
  ya que se pretende conservar el método físico.
  \item Sincronizar las notificaciones para los docentes en el
  sistema con el correo electrónico que proporciona el colegio.
  \item Mantener al día la información acerca de los exalumnos
  (si es que ellos lo desean) de la forma más accesible para ellos.
\end{itemize}
\ \newline
\normalsize{ \indent
Por otra parte, para describir el alcance que se presenta en esta
etapa se puede resumir en los dos primeros objetivos escritos en
esta parte. En otras palabras, las metas en este corto plazo son:
}
\begin{itemize}
  \item Establecer las condiciones que limitan el cambio de horarios
  por la jerarquía de importancia, es decir, cuáles son las que
  imposibilitan el movimiento, qué situaciones no tienen las
  validaciones suficientes para rotar (o que necesitan más horas en
  el movimiento cíclico) y cuándo hay derechos para la acción.
  \item Solucionar de forma correcta la situación de un relevamiento
  docente que a la vez debe tener cambios de horario. La causa de
  esto es que Coordinación Escolar controla los horarios, mientras
  que la rectora asigna los cargos.
  \item Crear un formulario digital que sirva para generar la 
  declaración jurada, así se imprime para su firma presencial y se
  almacenan el escaneo del mismo; junto con fotografías de los otros
  requisitos del legajo docente, los cuales pueden subir ellos mismos.
\end{itemize}
